\chapter{Les amines}

\textit{Les amines sont des composés organiques azotés dérivés de l'ammoniac par substitution d'un, deux ou trois atomes d'hydrogène par des groupes alkyles ou aryles. Elles jouent un rôle fondamental en chimie organique, tant comme bases que comme nucléophiles, et sont présentes dans de nombreuses molécules biologiques (alcaloïdes, acides aminés, neurotransmetteurs). Ce chapitre présente leur structure, leur nomenclature, leurs propriétés physiques et chimiques, ainsi que leurs méthodes de préparation.}

\section{Structure et classification des amines}

\subsection{Groupe amino}

Le groupe fonctionnel caractéristique des amines est le groupe \textbf{amino} : \ce{-NH2}, \ce{-NHR} ou \ce{-NR2}, selon le degré de substitution.

\begin{definition}[Groupe amino]
Le groupe amino est un groupement fonctionnel azoté dérivé de l'ammoniac \ce{NH3}, dans lequel un, deux ou trois atomes d'hydrogène sont remplacés par des groupes carbonés (alkyles ou aryles).
\end{definition}

\subsection{Classification des amines}

Les amines sont classées en trois catégories selon le nombre de groupes carbonés liés à l'atome d'azote :

\begin{itemize}
    \item \textbf{Amine primaire} : un seul groupe carboné lié à l'azote, formule générale \ce{R-NH2}.
    \item \textbf{Amine secondaire} : deux groupes carbonés liés à l'azote, formule générale \ce{R-NH-R$'$}.
    \item \textbf{Amine tertiaire} : trois groupes carbonés liés à l'azote, formule générale \ce{R-N(R$'$)-R$'$$'$}.
\end{itemize}

\begin{propriete}[Classification des amines]
Le nombre de groupes carbonés directement liés à l'atome d'azote détermine la classe de l'amine :
\begin{itemize}
    \item 1 groupe carboné : amine primaire (1$^\circ$)
    \item 2 groupes carbonés : amine secondaire (2$^\circ$)
    \item 3 groupes carbonés : amine tertiaire (3$^\circ$)
\end{itemize}
\end{propriete}

\begin{exemple}
\begin{itemize}
    \item \textbf{Amine primaire} : éthanamine \chemfig{CH_3-CH_2-NH_2}
    \item \textbf{Amine secondaire} : N-méthyléthanamine \chemfig{CH_3-CH_2-NH-CH_3}
    \item \textbf{Amine tertiaire} : N,N-diméthyléthanamine \chemfig{CH_3-CH_2-N(-CH_3)-CH_3}
\end{itemize}
\end{exemple}

\subsection{Comparaison avec l'ammoniac}

L'ammoniac \ce{NH3} peut être considéré comme l'amine la plus simple, où les trois atomes d'hydrogène sont liés à l'azote. Les amines sont des dérivés organiques de l'ammoniac.

\begin{remarque}
La classification des amines est analogue à celle des alcools, mais attention : un alcool tertiaire a le groupe \ce{-OH} lié à un carbone tertiaire, alors qu'une amine tertiaire a l'azote lié à trois carbones. Il ne faut pas confondre ces notions.
\end{remarque}

\section{Nomenclature des amines}

\subsection{Règles IUPAC}

La nomenclature systématique des amines suit les règles de l'IUPAC :

\begin{loi}[Nomenclature des amines]
\begin{itemize}
    \item On repère la chaîne carbonée la plus longue contenant le groupe \ce{-NH2} (pour les amines primaires).
    \item Le suffixe \textbf{-amine} remplace le \textbf{-e} final du nom de l'alcane correspondant.
    \item La position du groupe amino est indiquée par un numéro le plus petit possible.
    \item Pour les amines secondaires et tertiaires, on utilise le préfixe \textbf{N-} suivi du nom du substituant sur l'azote.
\end{itemize}
\end{loi}

\begin{exemple}
\begin{itemize}
    \item \chemfig{CH_3-NH_2} : méthanamine
    \item \chemfig{CH_3-CH_2-CH_2-NH_2} : propan-1-amine
    \item \chemfig{CH_3-CH(NH_2)-CH_3} : propan-2-amine
    \item \chemfig{CH_3-NH-CH_2-CH_3} : N-méthyléthanamine
    \item \chemfig{CH_3-N(-CH_3)-CH_2-CH_3} : N,N-diméthyléthanamine
\end{itemize}
\end{exemple}

\subsection{Noms courants}

Certaines amines portent des noms usuels encore employés :

\begin{itemize}
    \item \ce{CH3NH2} : méthylamine
    \item \ce{(CH3)2NH} : diméthylamine
    \item \ce{(CH3)3N} : triméthylamine
    \item \ce{C6H5NH2} : aniline (amine aromatique)
\end{itemize}

\begin{remarque}
L'aniline \chemfig{*6(-=-=-=)-NH_2} est une amine aromatique primaire très importante en chimie industrielle (colorants, médicaments).
\end{remarque}

\section{Propriétés physiques des amines}

\subsection{État physique et odeur}

Les amines aliphatiques de faible masse molaire sont des gaz ou des liquides volatils à température ambiante. Elles dégagent une odeur caractéristique, souvent désagréable (odeur de poisson pourri pour les petites amines).

\begin{experience}[Odeur des amines]
La triméthylamine \ce{(CH3)3N} est responsable de l'odeur caractéristique du poisson en décomposition. La putrescine \ce{H2N-(CH2)4-NH2} et la cadavérine \ce{H2N-(CH2)5-NH2} sont des diamines responsables des odeurs de putréfaction.
\end{experience}

\subsection{Point d'ébullition}

Les amines ont des points d'ébullition plus élevés que les alcanes de masse molaire comparable, mais généralement plus faibles que ceux des alcools correspondants.

\begin{propriete}[Points d'ébullition des amines]
\begin{itemize}
    \item Les amines primaires et secondaires peuvent former des liaisons hydrogène intermoléculaires (donneur et accepteur), ce qui élève leur point d'ébullition.
    \item Les amines tertiaires ne peuvent pas former de liaisons hydrogène entre elles (pas d'hydrogène sur l'azote), leur point d'ébullition est donc plus bas que celui des amines primaires ou secondaires de masse comparable.
    \item Les liaisons hydrogène N-H$\cdots$N sont moins fortes que les liaisons O-H$\cdots$O, donc les amines bouillent à température plus basse que les alcools.
\end{itemize}
\end{propriete}

\begin{exemple}
\begin{center}
\begin{tabular}{lccc}
\hline
Composé & Masse molaire (\si{\gram\per\mol}) & Point d'ébullition (\si{\celsius}) \\
\hline
Éthane \ce{CH3CH3} & 30 & -89 \\
Méthanamine \ce{CH3NH2} & 31 & -6 \\
Méthanol \ce{CH3OH} & 32 & 65 \\
\hline
\end{tabular}
\end{center}
\end{exemple}

\subsection{Solubilité dans l'eau}

Les petites amines (jusqu'à 5 ou 6 atomes de carbone) sont solubles dans l'eau grâce à la formation de liaisons hydrogène avec les molécules d'eau.

\begin{propriete}[Solubilité des amines]
\begin{itemize}
    \item Les amines primaires, secondaires et tertiaires de faible masse molaire sont solubles dans l'eau.
    \item La solubilité diminue lorsque la chaîne carbonée s'allonge (caractère hydrophobe du groupe alkyle).
    \item Les amines tertiaires sont généralement plus solubles que les amines primaires de même masse molaire (meilleure accessibilité du doublet de l'azote pour les liaisons hydrogène avec l'eau).
\end{itemize}
\end{propriete}

\section{Caractère basique des amines}

\subsection{Le couple ammonium/amine}

Les amines sont des bases de Brønsted : elles peuvent capter un proton \ce{H+} pour former un ion ammonium.

\begin{definition}[Ion ammonium]
L'ion ammonium est l'espèce obtenue par protonation d'une amine. Pour une amine primaire \ce{RNH2}, l'ion ammonium correspondant est \ce{RNH3+}.
\end{definition}

\begin{loi}[Couple acide-base ammonium/amine]
Le couple acide-base associé à une amine est :
\[\ce{RNH3+ <=> RNH2 + H+}\]
où \ce{RNH3+} est l'acide (ion ammonium) et \ce{RNH2} la base (amine).
\end{loi}

\subsection{Constante de basicité et pKa}

La basicité d'une amine se mesure par sa constante de basicité $K_b$ ou par le $pK_a$ du couple ammonium/amine correspondant.

\begin{propriete}[Constante de basicité]
Pour une amine \ce{RNH2} en solution aqueuse :
\[\ce{RNH2 + H2O <=> RNH3+ + HO-}\]
La constante de basicité s'écrit :
\[K_b = \frac{[\ce{RNH3+}][\ce{HO-}]}{[\ce{RNH2}]}\]
Le $pK_a$ du couple \ce{RNH3+/RNH2} est lié à $K_b$ par la relation :
\[pK_a + pK_b = 14 \quad \text{à } \SI{25}{\celsius}\]
\end{propriete}

\begin{aretenir}
Plus le $pK_a$ du couple ammonium/amine est élevé, plus l'amine est basique (plus elle capte facilement un proton).
\end{aretenir}

\subsection{Facteurs influençant la basicité}

La basicité des amines dépend de plusieurs facteurs :

\begin{propriete}[Ordre de basicité des amines]
\begin{itemize}
    \item \textbf{Effet inductif donneur} : les groupes alkyles sont donneurs d'électrons par effet inductif (+I), ce qui augmente la densité électronique sur l'azote et renforce le caractère basique.
    \item \textbf{Encombrement stérique} : dans les amines tertiaires, l'encombrement autour de l'azote peut gêner la solvatation de l'ion ammonium formé, ce qui diminue la basicité apparente.
    \item \textbf{Ordre général en solution aqueuse} : amine secondaire $>$ amine primaire $>$ amine tertiaire $>$ ammoniac.
\end{itemize}
\end{propriete}

\begin{exemple}
Valeurs de $pK_a$ des couples ammonium/amine correspondants (à \SI{25}{\celsius}) :
\begin{itemize}
    \item \ce{NH4+/NH3} : $pK_a = \num{9.2}$
    \item \ce{CH3NH3+/CH3NH2} : $pK_a = \num{10.6}$
    \item \ce{(CH3)2NH2+/(CH3)2NH} : $pK_a = \num{10.7}$
    \item \ce{(CH3)3NH+/(CH3)3N} : $pK_a = \num{9.8}$
\end{itemize}
\end{exemple}

\begin{remarque}
L'aniline \ce{C6H5NH2} est beaucoup moins basique que les amines aliphatiques ($pK_a \approx \num{4.6}$) car le doublet de l'azote est délocalisé dans le cycle aromatique (conjugaison avec le noyau benzénique), ce qui le rend moins disponible pour la protonation.
\end{remarque}

\subsection{Réaction avec les acides}

Les amines réagissent avec les acides pour former des sels d'ammonium.

\begin{experience}[Réaction d'une amine avec un acide]
L'éthanamine \ce{CH3CH2NH2} réagit avec l'acide chlorhydrique \ce{HCl} :
\[\ce{CH3CH2NH2 + HCl -> CH3CH2NH3+ + Cl-}\]
Le chlorure d'éthylammonium \ce{CH3CH2NH3+Cl-} est un sel soluble dans l'eau.
\end{experience}

\begin{loi}[Neutralisation des amines]
Les amines sont des bases qui réagissent avec les acides forts pour former des sels d'ammonium :
\[\ce{RNH2 + H3O+ -> RNH3+ + H2O}\]
\[\ce{RNH2 + HCl -> RNH3+Cl-}\]
\end{loi}

\subsection{Réaction avec l'eau}

Les amines en solution aqueuse sont des bases faibles qui réagissent partiellement avec l'eau.

\begin{exemple}
Solution de méthylamine dans l'eau :
\[\ce{CH3NH2 + H2O <=> CH3NH3+ + HO-}\]
Le pH d'une solution de méthylamine de concentration \SI{0.1}{\mol\per\litre} est d'environ \num{11.8} (basique).
\end{exemple}

\section{Nucléophilie des amines}

\subsection{Caractère nucléophile}

Le doublet non liant de l'atome d'azote confère aux amines un fort caractère nucléophile. Elles peuvent attaquer des sites électrophiles (carbones porteurs d'un atome d'halogène, carbonyles, etc.).

\begin{definition}[Nucléophilie des amines]
Une amine est un nucléophile : elle possède un doublet électronique libre sur l'azote qui peut attaquer un atome de carbone électrophile (déficitaire en électrons).
\end{definition}

\subsection{Réaction avec les halogénoalcanes : alkylation}

Les amines réagissent avec les halogénoalcanes par substitution nucléophile pour former des amines plus substituées.

\begin{experience}[Alkylation d'une amine]
L'ammoniac ou une amine primaire réagit avec un halogénoalcane :
\[\ce{CH3CH2Br + NH3 -> CH3CH2NH2 + HBr}\]
Cette réaction peut se poursuivre pour donner des amines secondaires, tertiaires, puis des sels d'ammonium quaternaire.
\end{experience}

\begin{propriete}[Alkylation successive]
L'alkylation des amines par les halogénoalcanes est une réaction qui peut conduire à un mélange de produits :
\[\ce{NH3 ->[RX] RNH2 ->[RX] R2NH ->[RX] R3N ->[RX] R4N+X-}\]
Le sel d'ammonium quaternaire \ce{R4N+X-} est le produit final de l'alkylation exhaustive.
\end{propriete}

\subsection{Réaction avec les dérivés carbonylés}

Les amines réagissent avec les aldéhydes et les cétones pour former des imines (bases de Schiff).

\begin{exemple}
Réaction de la méthanamine avec le méthanal :
\[\ce{CH3NH2 + HCHO -> CH3N=CH2 + H2O}\]
Il se forme une imine (base de Schiff).
\end{exemple}

\section{Préparation des amines}

\subsection{Alkylation de l'ammoniac}

C'est la méthode la plus directe : l'ammoniac réagit avec un halogénoalcane.

\begin{loi}[Alkylation de l'ammoniac]
\[\ce{R-X + NH3 -> RNH2 + HX}\]
Cette réaction produit un mélange d'amines primaire, secondaire, tertiaire et de sel d'ammonium quaternaire.
\end{loi}

\begin{remarque}
Pour favoriser la formation de l'amine primaire, on utilise un large excès d'ammoniac.
\end{remarque}

\subsection{Réduction des dérivés nitrés}

Les composés nitrés aromatiques peuvent être réduits en amines aromatiques.

\begin{experience}[Réduction du nitrobenzène]
Le nitrobenzène \chemfig{*6(-(-NO_2)=-=--=)} est réduit en aniline \chemfig{*6(-(-NH_2)=-=--=)} par action du fer en milieu acide ou par hydrogénation catalytique :
\[\ce{C6H5NO2 + 3H2 -> C6H5NH2 + 2H2O}\]
\end{experience}

\subsection{Réduction des nitriles et des amides}

Les nitriles et les amides peuvent être réduits en amines primaires.

\begin{exemple}
Réduction du propanenitrile en propan-1-amine :
\[\ce{CH3CH2CN + 2H2 -> CH3CH2CH2NH2}\]
\end{exemple}

\subsection{Réaction de Gabriel (synthèse d'amines primaires pures)}

La synthèse de Gabriel permet d'obtenir des amines primaires pures sans produits secondaires.

\begin{propriete}[Synthèse de Gabriel]
\begin{enumerate}
    \item Le phtalimide de potassium réagit avec un halogénoalcane pour former un N-alkylphtalimide.
    \item L'hydrolyse acide ou basique libère l'amine primaire pure.
\end{enumerate}
\end{propriete}

\section{L'essentiel à retenir}

\begin{synthese}
\subsection*{Les amines : ce qu'il faut savoir}

\begin{itemize}
    \item \textbf{Définition} : dérivés organiques de l'ammoniac \ce{NH3} par substitution d'hydrogènes par des groupes alkyles ou aryles.
    \item \textbf{Classification} : primaire (\ce{RNH2}), secondaire (\ce{R2NH}), tertiaire (\ce{R3N}).
    \item \textbf{Nomenclature} : suffixe \textbf{-amine}, position du groupe \ce{-NH2}, préfixe \textbf{N-} pour les substituants sur l'azote.
    \item \textbf{Propriétés physiques} : 
    \begin{itemize}
        \item Odeur caractéristique (poisson).
        \item Points d'ébullition : 1$^\circ$ > 2$^\circ$ > 3$^\circ$ (liaisons hydrogène).
        \item Solubilité dans l'eau pour les petites amines.
    \end{itemize}
    \item \textbf{Basicite} : 
    \begin{itemize}
        \item Couple \ce{RNH3+/RNH2}, $pK_a$ voisin de \numrange{9}{11} pour les amines aliphatiques.
        \item Ordre en solution aqueuse : 2$^\circ$ > 1$^\circ$ > 3$^\circ$ > \ce{NH3}.
        \item L'aniline est beaucoup moins basique ($pK_a \approx \num{4.6}$).
    \end{itemize}
    \item \textbf{Réactions} :
    \begin{itemize}
        \item Avec les acides : formation de sels d'ammonium.
        \item Avec l'eau : équilibre basique.
        \item Nucléophilie : alkylation par les halogénoalcanes.
        \item Avec les carbonylés : formation d'imines.
    \end{itemize}
    \item \textbf{Préparation} :
    \begin{itemize}
        \item Alkylation de l'ammoniac (mélange).
        \item Réduction des dérivés nitrés (amines aromatiques).
        \item Réduction des nitriles et amides.
        \item Synthèse de Gabriel (amines primaires pures).
    \end{itemize}
\end{itemize}

\subsection*{Pièges à éviter}
\begin{itemize}
    \item Ne pas confondre amine tertiaire et alcool tertiaire.
    \item L'ordre de basicité en solution aqueuse n'est pas celui attendu par le seul effet inductif (encombrement stérique).
    \item L'aniline est une base très faible à cause de la délocalisation du doublet de l'azote.
    \item L'alkylation de l'ammoniac donne toujours un mélange de produits.
\end{itemize}
\end{synthese}

\section{Méthodes \& savoir-faire}

\methode{Reconnaître la classe d'une amine}
Pour déterminer la classe d'une amine, compter le nombre d'atomes de carbone directement liés à l'atome d'azote.
\begin{itemize}
    \item 1 carbone lié à N : amine primaire.
    \item 2 carbones liés à N : amine secondaire.
    \item 3 carbones liés à N : amine tertiaire.
\end{itemize}

\begin{exemple}
Classer les amines suivantes :
\begin{itemize}
    \item \chemfig{CH_3-CH_2-NH-CH_3} : l'azote est lié à deux carbones (un éthyle et un méthyle) → amine secondaire.
    \item \chemfig{CH_3-CH(-NH_2)-CH_3} : l'azote est lié à un seul carbone → amine primaire.
    \item \chemfig{CH_3-N(-CH_3)-CH_2-CH_3} : l'azote est lié à trois carbones → amine tertiaire.
\end{itemize}
\end{exemple}

\methode{Nommer une amine selon l'IUPAC}
\begin{enumerate}
    \item Repérer la chaîne carbonée la plus longue contenant le groupe \ce{-NH2}.
    \item Numéroter la chaîne de façon à donner le plus petit indice au carbone portant le groupe amino.
    \item Nommer la chaîne principale avec le suffixe \textbf{-amine}.
    \item Pour les substituants sur l'azote, utiliser le préfixe \textbf{N-}.
\end{enumerate}

\begin{exemple}
Nommer \chemfig{CH_3-CH_2-CH(NH_2)-CH_3} :
\begin{itemize}
    \item Chaîne principale : 4 carbones → butane.
    \item Groupe \ce{-NH2} sur le carbone 2 → butan-2-amine.
\end{itemize}
Nommer \chemfig{CH_3-CH_2-NH-CH(CH_3)-CH_3} :
\begin{itemize}
    \item Chaîne principale : 3 carbones (propane) avec le groupe \ce{-NH-} en position 1.
    \item Substituant sur l'azote : méthyle → N-méthylpropan-1-amine.
\end{itemize}
\end{exemple}

\methode{Écrire l'équation de la réaction d'une amine avec un acide}
Une amine \ce{RNH2} réagit avec un acide fort \ce{HCl} ou \ce{H3O+} pour donner l'ion ammonium correspondant.

\begin{exemple}
Écrire l'équation de la réaction de la N-méthyléthanamine avec l'acide chlorhydrique.
\begin{itemize}
    \item Formule de l'amine : \ce{CH3CH2NHCH3} (amine secondaire).
    \item Réaction : \ce{CH3CH2NHCH3 + HCl -> CH3CH2NH2CH3+ + Cl-}.
    \item L'ion formé est le chlorure de N-méthyléthylammonium.
\end{itemize}
\end{exemple}

\methode{Calculer le pH d'une solution d'amine}
Pour une solution aqueuse d'amine de concentration $C$, on utilise la constante de basicité $K_b$ ou le $pK_a$ du couple ammonium/amine.

\begin{exemple}
Calculer le pH d'une solution de méthylamine \ce{CH3NH2} de concentration \SI{0.10}{\mol\per\litre} ($pK_a(\ce{CH3NH3+/CH3NH2}) = \num{10.6}$).
\begin{itemize}
    \item $pK_b = 14 - pK_a = 14 - 10,6 = 3,4$, donc $K_b = 10^{-3,4} = \num{3,98e-4}$.
    \item Équilibre : \ce{CH3NH2 + H2O <=> CH3NH3+ + HO-}.
    \item $K_b = \dfrac{x^2}{C - x} \approx \dfrac{x^2}{C}$ si $x \ll C$.
    \item $x = \sqrt{K_b \times C} = \sqrt{\num{3,98e-4} \times 0,10} = \sqrt{\num{3,98e-5}} = \num{6,31e-3}$.
    \item $[\ce{HO-}] = \num{6,31e-3}$ \si{\mol\per\litre}, $pOH = -\log(\num{6,31e-3}) = 2,2$.
    \item $pH = 14 - 2,2 = 11,8$.
\end{itemize}
\end{exemple}

\methode{Prévoir le produit d'une alkylation d'amine}
L'alkylation d'une amine par un halogénoalcane est une substitution nucléophile qui ajoute un groupe alkyle sur l'azote.

\begin{exemple}
Quel est le produit de la réaction entre l'éthanamine et le bromométhane en excès ?
\begin{itemize}
    \item L'éthanamine \ce{CH3CH2NH2} est une amine primaire.
    \item Avec un excès de \ce{CH3Br}, l'alkylation se poursuit jusqu'au sel d'ammonium quaternaire.
    \item Produit final : \ce{CH3CH2N+(CH3)3 Br-} (bromure d'éthyltriméthylammonium).
\end{itemize}
\end{exemple}

\methode{Distinguer une amine primaire, secondaire et tertiaire par réaction chimique}
Le test de Hinsberg permet de distinguer les trois classes d'amines.

\begin{exemple}
Le chlorure de benzènesulfonyle \ce{C6H5SO2Cl} réagit différemment :
\begin{itemize}
    \item Avec une amine primaire : formation d'un sulfonamide soluble en milieu basique.
    \item Avec une amine secondaire : formation d'un sulfonamide insoluble en milieu basique.
    \item Avec une amine tertiaire : pas de réaction (pas d'hydrogène sur l'azote).
\end{itemize}
\end{exemple}

\methode{Écrire le mécanisme de la réaction d'une amine avec un aldéhyde}
La réaction entre une amine primaire et un aldéhyde conduit à une imine (base de Schiff) par addition-élimination.

\begin{exemple}
Mécanisme de la réaction entre la méthanamine \ce{CH3NH2} et le méthanal \ce{HCHO} :
\begin{enumerate}
    \item Attaque nucléophile du doublet de l'azote sur le carbone carbonylé.
    \item Transfert de proton de l'azote vers l'oxygène.
    \item Élimination d'une molécule d'eau.
\end{enumerate}
Équation bilan : \ce{CH3NH2 + HCHO -> CH3N=CH2 + H2O}.
\end{exemple}

\methode{Préparer une amine primaire pure}
Pour obtenir une amine primaire sans produits secondaires, on utilise la synthèse de Gabriel.

\begin{exemple}
Préparation de la propan-1-amine par la méthode de Gabriel :
\begin{enumerate}
    \item Le phtalimide de potassium réagit avec le 1-bromopropane.
    \item Hydrolyse acide du N-propylphtalimide obtenu.
    \item Libération de la propan-1-amine pure.
\end{enumerate}
\end{exemple}